\begin{Exercise}[label={ex:flatten}]
    Si definisca una struttura dati per una lista concatenata in cui ogni nodo è 
    costituito da un valore intero, un puntatore al nodo successivo e un puntatore opzionale (\texttt{child}) 
    che identifica la testa di una sotto-lista gerarchica. 
    Si scriva un programma che costruisce tale struttura e implementa una funzione di inserimento e di stampa che gestisca
    la gerarchizzazione della struttura ed una funzione di "appiattimento".
    (flatten) che linearizza la gerarchia: ogni sotto-lista puntata da \texttt{child} deve essere inserita nella lista principale 
    immediatamente dopo il nodo genitore, modificando i puntatori esistenti senza allocare nuovi nodi. 
    Il programma visualizzi a video la sequenza lineare risultante.
    
    \textbf{Esempio:} data la lista gerarchica
    \begin{verbatim}
        1 -> 2 -> 5
             |
             3 -> 4
                  |
                  10
    \end{verbatim}
    la lista appiattita sarà:
    \begin{verbatim}
        1 -> 2 -> 3 -> 4 -> 10 -> 5
    \end{verbatim}
\end{Exercise}

\begin{Answer}[ref={ex:flatten}]
  	\lstinputlisting[style=cstylesmall]{src/flatten.c}
\end{Answer}